%% ma: L12-B16-0013
%% lop: 12
%% chuong: 5
%% bai: 16
%% dang: 12.16.10
%% loai: TN4
%% muc: TH
%% dap_an: "C"
%% nguon: "Ảnh giáo viên gửi 04/11/2026 (Huỳnh Văn Ánh, Chuyên đề V - Phương pháp tọa độ trong không gian), Đề test theo bài, Phần I Câu 4"
%% kien_thuc: "Bài 16 (góc giữa đường thẳng và mặt phẳng), Bài 15 (vectơ chỉ phương từ phương trình chính tắc), Bài 14 (vectơ pháp tuyến)"
%% trang_thai: cho-duyet
%% ly_do: "Dạng 12.16.10 mới đề xuất, chờ giáo viên duyệt."
%% ghi_chu: "Đáp án gốc C khớp (sympy). Câu hỏi côsin nên cần thêm bước đổi từ sin sang cos; phân vân giữa TH và VD, chọn mức thấp hơn."
\begin{ex}
Trong không gian với hệ trục tọa độ $Oxyz$, cho đường thẳng $d\colon\dfrac{x-2}{2}=\dfrac{y+1}{-1}=\dfrac{z-1}{2}$ và mặt phẳng $(P)\colon 2x-y-z-5=0$. Hãy tính côsin góc tạo bởi đường thẳng $(d)$ và mặt phẳng $(P)$.
\choice{$\dfrac1{\sqrt6}$}{$\dfrac12$}{\True $\dfrac{\sqrt{30}}{6}$}{$\dfrac{\sqrt3}{2}$}
\loigiai{$d$ có vectơ chỉ phương $\vec u=(2;-1;2)$, $(P)$ có vectơ pháp tuyến $\vec n=(2;-1;-1)$.

Gọi $\alpha$ là góc giữa $d$ và $(P)$ ($0^\circ\le\alpha\le90^\circ$): $\sin\alpha=\dfrac{|\vec u\cdot\vec n|}{|\vec u|\cdot|\vec n|}=\dfrac{|4+1-2|}{3\cdot\sqrt6}=\dfrac1{\sqrt6}$.

Suy ra $\cos\alpha=\sqrt{1-\sin^2\alpha}=\sqrt{1-\dfrac16}=\dfrac{\sqrt{30}}{6}$. Chọn C.}
\end{ex}
